% GENERATED FILE. Edit canonical lessons, then run npm run generate:books.
% canonical-source-hash: fnv1a64:6af062c2b78d6fe5
% canonical-lessons: PT-C395-ganhar, PT-C395-lucro, PT-C395-perda, PT-C395-falir, PT-C395-endividado

\chapter{Ganhos e perdas}
\label{ch:pt-ganhos-e-perdas}
\hlchaptermodality{395}

\begin{chapteropening}
\textbf{By the end of this chapter:} \emph{I can talk about earning, profit and loss, going bankrupt and being in debt.}

Complete the last lesson of chapter 395: I can talk about earning, profit and loss, going bankrupt and being in debt.
\end{chapteropening}

\section[ganhar]{ganhar — to earn; to win; to gain}

\label{lesson:PT-C395-ganhar}



\begin{quote}
Before the new one: say the Portuguese for to refund, then the Portuguese for the credit.
\end{quote}

\subsection*{You'll want to know: ganhar}
\textbf{ganhar} — \textquotedblleft{}to earn; to win; to gain\textquotedblright{}.

\begin{quote}
Ela ganha mil e quinhentos euros por mês.
\end{quote}

\begin{quote}
\emph{She earns one thousand five hundred euros a month.}
\end{quote}

\emph{Ganhar a vida} — to earn a living; \emph{ganhar um jogo} — to win a game.

\begin{cousinweb}
\emph{ganhar} is usually traced to a Germanic verb meaning \textquotedblleft{}to graze, to hunt for food\textquotedblright{}, through Vulgar Latin \emph{wadaniāre} (EtymDB); French \emph{gagner} and Italian \emph{guadagnare} are its cousins.
\end{cousinweb}

\begin{grammarlens}[title={What you've built}]
One more word for telling, explaining and describing at B1.
\end{grammarlens}

\subsection*{Your turn}
\begin{itemize}
\raggedright
  \item \emph{Say it:} \emph{ganhar}
  \item \emph{Say it:} \emph{ganhar}, once more
  \item \emph{Say it:} a sentence of your own with \emph{ganhar}
  \item \emph{Recall:} say the Portuguese for to refund, then the Portuguese for the credit, then say \emph{ganhar} again
\end{itemize}

\begin{tcolorbox}[breakable,colback=teal!4,colframe=teal!35!black,title={Before you move on}]
What does ganhar mean? (\textquotedblleft{}To earn; to win; to gain\textquotedblright{}.) Say it once more.
\end{tcolorbox}
\section[o lucro]{o lucro — the profit}

\label{lesson:PT-C395-lucro}



\begin{quote}
Before the new one: say the Portuguese for the credit, then the Portuguese for to earn.
\end{quote}

\subsection*{You'll want to know: o lucro}
\textbf{o lucro} — \textquotedblleft{}the profit\textquotedblright{}.

\begin{quote}
A loja teve lucro pela primeira vez.
\end{quote}

\begin{quote}
\emph{The shop made a profit for the first time.}
\end{quote}

\emph{Dar lucro} — to make a profit.

\begin{cousinweb}
\emph{lucro} is a learned borrowing from Latin \emph{lucrum} \textquotedblleft{}gain, profit\textquotedblright{} (EtymDB, Lewis \& Short).
\end{cousinweb}

\begin{grammarlens}[title={What you've built}]
One more word for telling, explaining and describing at B1.
\end{grammarlens}

\subsection*{Your turn}
\begin{itemize}
\raggedright
  \item \emph{Say it:} \emph{o lucro}
  \item \emph{Say it:} \emph{o lucro}, once more
  \item \emph{Say it:} a sentence of your own with \emph{o lucro}
  \item \emph{Recall:} say the Portuguese for the credit, then the Portuguese for to earn, then say \emph{o lucro} again
\end{itemize}

\begin{tcolorbox}[breakable,colback=teal!4,colframe=teal!35!black,title={Before you move on}]
What does o lucro mean? (\textquotedblleft{}The profit\textquotedblright{}.) Say it once more.
\end{tcolorbox}
\section[a perda]{a perda — the loss}

\label{lesson:PT-C395-perda}



\begin{quote}
Before the new one: say the Portuguese for to earn, then the Portuguese for the profit.
\end{quote}

\subsection*{You'll want to know: a perda}
\textbf{a perda} — \textquotedblleft{}the loss\textquotedblright{}.

\begin{quote}
A empresa teve grandes perdas este ano.
\end{quote}

\begin{quote}
\emph{The company had big losses this year.}
\end{quote}

\emph{Perda de tempo} — waste of time.

\begin{cousinweb}
\emph{perda} is inherited from Vulgar Latin \emph{perdita} \textquotedblleft{}a thing lost\textquotedblright{} (EtymDB), from \emph{perdere} \textquotedblleft{}to destroy, to lose\textquotedblright{} (Lewis \& Short).
\end{cousinweb}

\begin{grammarlens}[title={What you've built}]
One more word for telling, explaining and describing at B1.
\end{grammarlens}

\subsection*{Your turn}
\begin{itemize}
\raggedright
  \item \emph{Say it:} \emph{a perda}
  \item \emph{Say it:} \emph{a perda}, once more
  \item \emph{Say it:} a sentence of your own with \emph{a perda}
  \item \emph{Recall:} say the Portuguese for to earn, then the Portuguese for the profit, then say \emph{a perda} again
\end{itemize}

\begin{tcolorbox}[breakable,colback=teal!4,colframe=teal!35!black,title={Before you move on}]
What does a perda mean? (\textquotedblleft{}The loss\textquotedblright{}.) Say it once more.
\end{tcolorbox}
\section[falir]{falir — to go bankrupt}

\label{lesson:PT-C395-falir}



\begin{quote}
Before the new one: say the Portuguese for the profit, then the Portuguese for the loss.
\end{quote}

\subsection*{You'll want to know: falir}
\textbf{falir} — \textquotedblleft{}to go bankrupt\textquotedblright{}.

\begin{quote}
A loja do meu tio faliu.
\end{quote}

\begin{quote}
\emph{My uncle's shop went bankrupt.}
\end{quote}

Defective verb: forms like \emph{faliu} and \emph{falimos} are common. \emph{A falência} — bankruptcy.

\begin{cousinweb}
\emph{falir} is usually traced to Italian \emph{fallire} \textquotedblleft{}to go bankrupt\textquotedblright{}, which EtymDB derives from Latin \emph{fallere} (EtymDB), \textquotedblleft{}to deceive, to disappoint\textquotedblright{} (Lewis \& Short).
\end{cousinweb}

\begin{grammarlens}[title={What you've built}]
One more word for telling, explaining and describing at B1.
\end{grammarlens}

\subsection*{Your turn}
\begin{itemize}
\raggedright
  \item \emph{Say it:} \emph{falir}
  \item \emph{Say it:} \emph{falir}, once more
  \item \emph{Say it:} a sentence of your own with \emph{falir}
  \item \emph{Recall:} say the Portuguese for the profit, then the Portuguese for the loss, then say \emph{falir} again
\end{itemize}

\begin{tcolorbox}[breakable,colback=teal!4,colframe=teal!35!black,title={Before you move on}]
What does falir mean? (\textquotedblleft{}To go bankrupt\textquotedblright{}.) Say it once more.
\end{tcolorbox}
\section[endividado]{endividado — in debt}

\label{lesson:PT-C395-endividado}



\begin{quote}
Before the new one: say the Portuguese for the loss, then the Portuguese for to go bankrupt.
\end{quote}

\subsection*{You'll want to know: endividado}
\textbf{endividado} — \textquotedblleft{}in debt\textquotedblright{}.

\begin{quote}
Muitas famílias estão endividadas.
\end{quote}

\begin{quote}
\emph{Many families are in debt.}
\end{quote}

Built on \emph{dívida} — debt.

\begin{cousinweb}
\emph{endividado} is built from \emph{en-} and \emph{dívida} \textquotedblleft{}debt\textquotedblright{}, from Latin \emph{dēbita} \textquotedblleft{}things owed\textquotedblright{}, on \emph{dēbēre} \textquotedblleft{}to owe\textquotedblright{} (Lewis \& Short).
\end{cousinweb}

\begin{grammarlens}[title={What you've built}]
One more word for telling, explaining and describing at B1.
\end{grammarlens}

\subsection*{Your turn}
\begin{itemize}
\raggedright
  \item \emph{Say it:} \emph{endividado}
  \item \emph{Say it:} \emph{endividado}, once more
  \item \emph{Say it:} a sentence of your own with \emph{endividado}
  \item \emph{Recall:} say the Portuguese for the loss, then the Portuguese for to go bankrupt, then say \emph{endividado} again
\end{itemize}

\begin{tcolorbox}[breakable,colback=teal!4,colframe=teal!35!black,title={Before you move on}]
What does endividado mean? (\textquotedblleft{}In debt\textquotedblright{}.) Say it once more.
\end{tcolorbox}
